\chapter{Introduction}

Mathematics is permeated by this notion of what one
would call an ``ideal'' object. This notion of ``ideal'' often refers to the
simplest object that meets a specific criterion, such as being the smallest
element in a set, or the one through which all others factor uniquely.
This particular description of ``ideal'' reveals an interesting
interconnection between categories and optimization \cite{Khadh2026bridging}.
Many of the constructions mathematics singles
out as canonical are optimal in this same sense: they are the best object of a
collection, with respect to an ordering or universal property carried by that
collection. This thesis aims to explore this idea and provide a formal and
intuitive view on such topics.

Consider for instance the familiar transitivity condition for a relation
$R\subseteq X\times X$:
\[
\bigl(\exists y\in X:\;xRy\wedge yRz\bigr)\implies xRz,
\qquad x,z\in X.
\]

Here the premise demands the existence of an intermediate object, namely $y$.
Nothing forces $y$ to be unique; in fact, it is often the case that we have
a set of such objects. The connection with an ``ideal'' object can be made
explicit if we change the settings of the problem. Consider the indicator function
representation of $R$ in the sense of convex analysis \cite{rockafellar1970convex},
$R:X\times X\to\{0,\infty\}$, defined by
\[
R(x,y)=
\begin{cases}
0, & \text{if }xRy,\\
\infty, & \text{otherwise}.
\end{cases}
\]
Then the above is the same as:
\[
\inf_{y\in X}\bigl(R(x,y)+R(y,z)\bigr)=0
\implies R(x,z)=0,
\]
where $\inf\varnothing=\infty$.
Since the values are either zero or infinity, this is also equivalent to
\[
\inf_{y\in X}\bigl(R(x,y)+R(y,z)\bigr) 
\ge R(x,z)
\]

The resemblance is not mere coincidence: these are particular cases of
a more general construction using quantales
\cite{lawvere1973metric}.
More generally, let $Q$ be a unital quantale: a complete lattice equipped
with an associative operation $\otimes$, with unit $e$, that preserves
arbitrary joins in each argument.
By interpreting $\exists$ and $\inf$ through the lattice join
operator $\bigvee$, and $\wedge$ and $+$ through the monoidal operator
$\otimes$, we obtain, for a $Q$-valued relation
$R:X\times X\to Q$, the generalized transitivity condition
\[
\bigvee_{y\in X}R(x,y)\otimes R(y,z)\le_Q R(x,z).
\]

Notice that for the cost quantale $([0,\infty],\ge,+)$ we have
$a\le_Q b \iff a\ge b$, so this replacement recovers exactly the same
expression as above.

However, there is still more to it: as one may notice, the commonalities
between these equations go beyond the simply visual structure;
they actually share an \textit{algebraic} structure.
Explicitly, all the expressions on the left of the abstract inequality
perform the very important role of \textit{variable elimination}.
The two terms share the same middle variable $y$, and taking the join
over $y$ produces a new expression depending only on $x$ and $z$.

A join, however, already has an optimization characterization:
for $A\subseteq Q$, $\bigvee A$ is the least of the upper bounds of $A$,
the ``ideal'' upper bound inside the order of $Q$.
Here the collection being aggregated is
\[
A_{x,z}=\{R(x,y)\otimes R(y,z):y\in X\}.
\]
Eliminating the intermediate variable therefore computes an optimal
aggregate value. The choice of $Q$ fixes both the order used for aggregation and
the operation used to combine values.
Ordinary existential quantification is thus a particular case of this
order-theoretic aggregation, and its indicator function makes the
connection with optimization exact.

This variable elimination is precisely how composition of
quantale-valued relations is defined. Given
$R:X\times Y\to Q$ and $S:Y\times Z\to Q$, we have
\[
(S\circ R)(x,z)
:=
\bigvee_{y\in Y}R(x,y)\otimes S(y,z).
\]
The generalized transitivity condition above is therefore simply
\[
R\circ R\le_Q R,
\]
with the inequality understood pointwise.
In this setting, composition is variable elimination, and the join
gives it the optimization interpretation described above.
Sections~\ref{sec:quantale-relations}
and~\ref{sec:composability} make this precise.

Category theory gives us an abstraction to describe such operations:
it explicitly encodes the notion of composition inside its own definition.
Quantale-valued relations form a category under the operation above,
with identity taking value $e$ on the diagonal and $\bot$ elsewhere.
More generally, composition in a category $\mathcal C$ is given by
\[
\circ:\mathcal C(Y,Z)\times\mathcal C(X,Y)
\longrightarrow\mathcal C(X,Z),
\]
subject to associativity and identity laws.
Its particular expression as a join over an intermediate variable
uses the additional quantale-valued structure.

We show that, in the categories treated here, solving an optimization
problem amounts to composing morphisms, and that composition carries out the
variable elimination such a solution requires. Convex optimization then admits a
single compositional presentation \cite{stein2024compositional, hanks2024compositionalopt}, 
in which optimization is an algebraic operation rather than a technique applied from outside.
That is, optimization and variable elimination are a syntactic operation,
namely the composition $\circ$ of two objects $A$ and $B$ that share a common variable, or type.

This compositional view of optimization becomes especially apparent when analyzed through
the language of string diagrams \cite{joyal1991geometry, selinger2010survey, piedeleu2023introductionstringdiagramscomputer}.
Originating as a rigorous graphical language for
monoidal categories, string diagrams provide an intuitive yet mathematically
precise syntax in which the sequential and parallel composition of morphisms
correspond directly to vertical and horizontal juxtaposition of diagrams
(e.g., representing computational processes)
\cite{joyal1991geometry, maclane1998categories}.
Graphical languages are remarkably well suited to the case at hand.
They make the semantics of the underlying mathematics explicit at the syntax level,
in the sense that one doesn`t need to go through atomic and subtle operations in order
to express the core ideas of a theory,
\textit{but the system's main properties are explicit in the syntax operations themselves}.

For example, consider the following diagram. Because of the monoidal structure
one may parse it by reading in two manners: composing the boxes in the same line
and then stack them togheter, or stacking the boxes that are vertically aligned 
and then composing then sequentially, as the dotted boxes shows.

\begin{figure}[H]
    \begin{center}
\begin{tikzpicture}[axpic]
  \pic[val=A] (g1) at (0.500,0.500) {diagram};
  \begin{scope}[shift={(1.500,0.000)}]
  \pic[val=B] (g2) at (0.500,0.500) {diagram};
  \end{scope}
  \pic[val=C] (g3) at (0.500,-0.500) {diagram};
  \begin{scope}[shift={(1.500,0.000)}]
  \pic[val=D] (g4) at (0.500,-0.500) {diagram};
  \end{scope}
  \draw[dotted, thick] (0,-0.1) rectangle (2.5,-1);
  \draw[dotted, thick] (0,0.1) rectangle (2.5,1);
  \draw[dotted, thick] (-0.2,-1.2) rectangle (2.7,1.2);

  \draw (g1-out-0) to[out=0,in=180] (g2-in-0);
  \draw (g3-out-0) to[out=0,in=180] (g4-in-0);
\end{tikzpicture}
\begin{tikzpicture}[axpic]
  \node (rel4) at (3.000,0.500) {$=$};
\end{tikzpicture}
\begin{tikzpicture}[axpic]
  \pic[val=A] (g1) at (0.500,0.500) {diagram};
  \begin{scope}[shift={(1.500,0.000)}]
  \pic[val=B] (g2) at (0.500,0.500) {diagram};
  \end{scope}
  \pic[val=C] (g3) at (0.500,-0.500) {diagram};
  \begin{scope}[shift={(1.500,0.000)}]
  \pic[val=D] (g4) at (0.500,-0.500) {diagram};
  \end{scope}
  \draw[dotted, thick] (0,-1.0) rectangle (1,1.0);
  \draw[dotted, thick] (2.5,-1.0) rectangle (1.5,1.0);
  \draw[dotted, thick] (-0.2,-1.2) rectangle (2.7,1.2);
  \draw (g1-out-0) to[out=0,in=180] (g2-in-0);
  \draw (g3-out-0) to[out=0,in=180] (g4-in-0);
\end{tikzpicture}
\end{center}
\end{figure}

This denotes the fact that the composition and the monoidal product
interact nicely, or, formally, that the operator $\otimes$ is actually
an endofunctor, so $a;b \otimes c;d = a \otimes c; b \otimes d$. Observe how
this entire property is underloaded by the syntax, so one don't even need to keep
track of it but simply parse the diagram in any direction. This is what we mean
when we say that the main properties are in the syntax itself, we don't get to
invoke the definition or property to show the equality,
but the very depicting (or drawing) embeds the property.

A fine example, that will be later explained, is one that uses the very syntax
developed in this work. Consider the equivalence between these two Linear Programs
\[
f(y) = \inf_{x_1,x_2} \{ x_1 + x_2 \mid \; a_1 x_1 + a_2 x_2 = y\}
\]
\[
f_2(y) = \inf_{w_1,w_2} \{ inf_{x_1} \{ x_1 \mid a_1 x_1 = w_1 \} + inf_{x_2} \{ x_2 \mid a_2 x_2 = w_2 \} \mid w_1 + w_2 = y \}
\]

These two functions are, in fact, the same. To see this, introduce the change of variables
$ w_1=a_1x_1,\; w_2=a_2x_2. $

Assuming \(a_1,a_2\neq 0\), each inner optimization problem has a unique feasible point,
so substituting \(w_i=a_ix_i\) yields
$ f_2(y) = \inf_{a_1x_1+a_2x_2=y} \left\{ x_1+x_2 \right\} = f(y). $

Thus, the second formulation does not change the underlying optimization problem;
rather, it decomposes the original constraint through the intermediate variables
\(w_1\) and \(w_2\). Now look at the diagram denoting these same problems.

\begin{figure}[H]
    \begin{center}
\begin{tikzpicture}[axpic]
  \pic (g2) at (-0.500,0.500) {linear};
  \pic[val=a_1] (g1) at (0.500,0.500) {scalar};
  \pic (g2) at (-0.500,-0.500) {linear};
  \pic[val=a_2] (g3) at (0.500,-0.500) {scalar};
  \pic (g4) at (2.,0.0) {multiplication};
  \draw[dotted, thick] (0,-0.9) rectangle (2.6,1);
  \draw[dotted, thick] (-0.9,-0.9) rectangle (-0.1,1);
  \draw[dotted, thick] (-1.0,-1.2) rectangle (2.7,1.2);

  \draw (g1-out-0) to[out=0,in=180] (g4-in-0);
  \draw (g3-out-0) to[out=0,in=180] (g4-in-1);
\end{tikzpicture}
\begin{tikzpicture}[axpic]
  \node (rel4) at (3.000,0.500) {$=$};
\end{tikzpicture}
\begin{tikzpicture}[axpic]
  \pic (g2) at (-0.500,0.500) {linear};
  \pic[val=a_1] (g1) at (0.500,0.500) {scalar};
  \pic (g2) at (-0.500,-0.500) {linear};
  \pic[val=a_2] (g3) at (0.500,-0.500) {scalar};
  \pic (g4) at (2.,0.0) {multiplication};
  \draw[dotted, thick] (-.9,0.1) rectangle (1.0,1.1);
  \draw[dotted, thick] (-.9,-0.1) rectangle (1.0,-.9);
  \draw[dotted, thick] (1.1,-0.9) rectangle (2.6,1.1);
  \draw[dotted, thick] (-1.0,-1.2) rectangle (2.7,1.2);

  \draw (g1-out-0) to[out=0,in=180] (g4-in-0);
  \draw (g3-out-0) to[out=0,in=180] (g4-in-1);
\end{tikzpicture}
\end{center}
\end{figure}

Once again, both problems are depicted as the very same diagram, being just
a different order of parsing it. In this case all the variable changing was 
hidden by the syntax, and only the main semmantics -- the \textit{problem itself} --
is exposed.

Other areas have far off recognize the benefits of such denotation, mainly in 
engineering fields such as \emph{quantum computing} \cite{coecke2011interacting} and
\emph{electronics} \cite{baez2018compositional}.
More recently, mathematicians and computer scientists have also 
built upon this foundation deriving a rich family of \emph{diagrammatic} 
or \emph{graphical calculi},
offering complete axiomatisations of various algebraic structures purely through
equational diagram rewriting \cite{selinger2010survey, bonchi2017interacting}.
These calculi have proved remarkably versatile, finding applications across a
wide range of fields:  \emph{linear algebra} \cite{bonchi2017interacting, bonchi2019affine, defreitas2025generalizinginvertiblematrixtheorem};
and in \emph{optimisation and probability} \cite{stein2024gqa, stein2025compositional, bonchi2021polyhedral}. 

Throughout this work we explore how, step by step,
one can enrich the already known theories for linear
\cite{bonchi2017interacting} algebra by adding one new generator at a time,
along with a particular
set of axioms, and derive even richer structures for convex problems \cite{stein2024compositional}.

The development followed a strict linear discipline, adding exactly one new
generator and verifying that the resulting algebra is
sound, expressive, and complete with respect to a precisely defined semantic
category. Each new class of optimization problems then costs a single
additional generator and a finite set of axioms.
This is made precise by the normal form and completeness results of
Theorem~\ref{thm:normalformgpa} and Theorem~\ref{thm:NormalFormGLP}, which follow the
same pattern already established for graphical linear and affine algebra
\cite{bonchi2017interacting, bonchi2019affine} and for their quadratic
\cite{stein2024gqa} and polyhedral \cite{bonchi2021polyhedral} extensions.

We start by recalling how the Symmetric Monoidal Theory of Graphical Affine Algebra (\textbf{GAA}) \cite{bonchi2019affine}
can be extended into the even more general theory of Symmetric Inequality Monoidal Theory.
This is achieved by adding the inequality generator $\inlinegen{ge}$. Geometrically,
this allows our algebra to go beyond affine spaces and start representing polyhedral figures as linear inequality constraints in optimization.
\[
	\{| x \in P\} =
	\begin{cases}
		0, & Ax+b\ge 0,\\
		+\infty, & \text{otherwise}.
	\end{cases}
\]

Inspired by these observations, and parallel works in this direction \cite{stein2024gqa, lawvere1973metric},
we make our central contribution in the final chapter: the construction and completeness of
a graphical algebra for parametric linear programming. In this
setting, a linear program is represented as a morphism that
can be composed with other optimization blocks, leading to the idea of
\textit{composing optimization problems}. This gives a symbolic language
in which the elimination of internal variables
corresponds directly to categorical composition, allowing optimization problems
to be manipulated diagrammatically before computation.

In particular, every piecewise affine convex function can be represented as
the value function of a linear program \cite{rockafellar1970convex},
which means the Graphical Linear Programming \textbf{GLP}
theory provides a syntax for a broad class of nonsmooth convex models that appear in optimization,
control, economics, and machine learning.
The diagrammatic syntax therefore erases the distinction between expressing a
linear program as an optimization problem and expressing it as a piecewise
affine convex function. As the example above, we have only one syntax to denote
one object, instead of two distinct ways to syntactically depict it.

\paragraph{Thesis overview}
In Section~\ref{sec:quantale-relations} we set up the algebraic and categorical
setting on which the rest of the work is built, introducing quantales and
relations and the interchangeable role they play in a categorical view. Next, in
Section~\ref{sec:composability}, we apply this theory to convex optimization
\cite{stein2024compositional, Khadh2026bridging}, showing how several of its
central constructs fall within the same categorical framework.

Chapter~\ref{sec:gpa} presents the already-known graphical theory of Polyhedral Algebra
\cite{bonchi2021polyhedral}, where some of the known proofs are recalled and refined.

Finally, in Chapter~\ref{sec:glp}, we present the central contribution of this thesis:
an entirely new Graphical Algebra, \textbf{GLP}, dedicated to expressing parametric Linear Programs.
Building on the polyhedral generators of Chapter~\ref{sec:gpa}, we introduce a single new generator
for the linear objective function and derive a complete axiomatization for right-hand side parametric
linear programming, proving that \textbf{GLP} presents the prop $\mathbf{LPRel}$.
This work represents a first step toward a complete diagrammatic calculus for convex optimization.
